\section*{\zihao{2} \centering 摘要}

\vskip 0.5cm

本研究从血管网络的几何结构出发，通过数学建模推导基础代谢率与体重之间的幂律关系。
基于对数学生理学的分析，我们建立了血管分支网络模型，推导出基础代谢率与体重的3/4次幂成正比，即著名的Kleiber定律。
为验证这一理论推断，我们收集了多种哺乳动物的生理数据（包括体重、静息心率和基础代谢率），
利用Python进行非线性拟合和统计分析。拟合结果表明，BMR幂指数为0.7705，与理论值0.75仅偏差2.73\%；
心率幂指数为-0.2620，与理论值-0.25偏差4.79\%。决定系数$R^2 > 0.99$，真实数据与Kleiber定律高度吻合，
验证了我们推论的正确性。

\textbf{关键词:} Kleiber定律\quad 基础代谢率\quad 幂律关系\quad 心率\quad 异速生长\quad 数学生理学
\addcontentsline{toc}{section}{摘要}

\clearpage
\section*{\zihao{2} \centering \textbf{Abstract} }

This study investigates the power-law relationship between basal metabolic rate (BMR) and body mass through 
mathematical modeling based on the geometric structure of vascular networks. 
By analyzing the fractal properties of branching vasculature, we establish a vascular branching network model 
that theoretically derives BMR proportional to $M^{3/4}$, namely the famous Kleiber's Law. 
To validate this theoretical prediction, we collected physiological data from multiple mammalian species 
(including body mass, resting heart rate, and BMR), and performed nonlinear fitting and statistical analysis using Python. 
The fitting results show that the BMR scaling exponent is 0.7705, deviating only 2.73\% from the theoretical value of 0.75; 
the heart rate scaling exponent is $-$0.2620, with a 4.79\% deviation from the theoretical value of $-$0.25. 
The coefficient of determination $R^2 > 0.99$ indicates that real data highly agrees with Kleiber's Law, 
confirming the correctness of our derivation.

\vskip 0.2cm
\noindent\textbf{Keywords:} Kleiber's Law\quad Basal Metabolic Rate\quad Power-law Scaling\quad Heart Rate\quad Allometric Growth\quad Mathematical Physiology
\addcontentsline{toc}{section}{Abstract}
